% ===========================================================================
\chapter{Classificação Tecnológica de Patentes (IPC)}
\label{cap_classificacao_ipc}
% ===========================================================================

A Classificação Internacional de Patentes (\textit{International Patent Classification} -- IPC), administrada pela Organização Mundial da Propriedade Intelectual \cite{wipo2023}, é um sistema hierárquico universal que indexa documentos patentários segundo as áreas técnicas a que pertencem as invenções.

Considerando o escopo inventivo, as rotinas computacionais e a finalidade prática do software, o sistema \textbf{SIGEA-GTT} classifica-se estrategicamente na interseção de duas classes fundamentais da tabela IPC: a classe médica e de saúde digital (\textbf{G16H}) e a classe de sistemas educacionais e de treinamento simulado (\textbf{G09B}).

% ===========================================================================
\section{Classe G16H: TIC Aplicada à Saúde}
\label{sec_classe_g16h}
% ===========================================================================

A classe \textbf{G16H} contempla invenções implementadas em computador voltadas à gestão, diagnóstico, monitoramento e análise de dados clínicos hospitalares. No SIGEA-GTT, o enquadramento dá-se especificamente nas seguintes subclasses e grupos:

\begin{enumerate}
    \item \textbf{G16H 40/00 (TIC especialmente adaptada para a gestão de instalações médicas ou empresas de saúde)}:
    \begin{itemize}
        \item \textit{G16H 40/20}: Tecnologias voltadas à gestão ou administração de cuidados de saúde, clínicas ou hospitais;
        \item \textit{Aplicação no SIGEA-GTT}: O módulo de governança acadêmica de turmas, o cadastro canônico de disciplinas e a emissão de relatórios gerenciais e taxas hospitalares oficiais do IHI (taxa de eventos por 1.000 pacientes-dia e taxa por 100 admissões) enquadram-se nesse segmento.
    \end{itemize}

    \item \textbf{G16H 50/00 (TIC especialmente adaptada para diagnósticos médicos, vigilância médica ou simulação médica)}:
    \begin{itemize}
        \item \textit{G16H 50/30}: Tecnologias direcionadas à avaliação ou cálculo de riscos à saúde do paciente;
        \item \textit{Aplicação no SIGEA-GTT}: O motor determinístico de indexação de gatilhos clínicos nos prontuários simulados, a classificação de severidade de dano conforme a taxonomia do NCC MERP (categorias A a I) e a aplicação dos eixos de gravidade, urgência e tendência da Matriz GUT operam sob esta classe.
    \end{itemize}
\end{enumerate}

% ===========================================================================
\section{Classe G09B: Sistemas e Instrumentos de Ensino e Treinamento}
\label{sec_classe_g09b}
% ===========================================================================

A classe \textbf{G09B} engloba dispositivos, matrizes lógicas e plataformas computacionais concebidos para a educação profissional, instrução assistida por computador e simulação de cenários operacionais. No SIGEA-GTT, destacam-se os seguintes grupos:

\begin{enumerate}
    \item \textbf{G09B 9/00 (Sistemas simuladores para ensino ou treinamento)}:
    \begin{itemize}
        \item \textit{Aplicação no SIGEA-GTT}: A arquitetura de \textit{sandbox} clínica do SIGEA-GTT, que projeta prontuários sintéticos de alta fidelidade persistidos em colunas semiestruturadas \texttt{JSONB} sob o controle de um cronômetro regressivo com alertas suaves de tempo aos 15 e 18 minutos;
    \end{itemize}

    \item \textbf{G09B 19/00 e G09B 23/28 (Ensino não previsto em outros grupos e modelos para ensino médico)}:
    \begin{itemize}
        \item \textit{Aplicação no SIGEA-GTT}: O espelho de correção pedagógico automatizado, que justapõe as marcações discentes ao gabarito docente oficial, e o campo obrigatório de parecer formativo docente concebido para o desenvolvimento do raciocínio causal nos futuros profissionais de saúde.
    \end{itemize}
\end{enumerate}

% ===========================================================================
\section{A Convergência Tecnológica como Núcleo da Invenção}
\label{sec_convergencia_tecnologica}
% ===========================================================================

A maioria absoluta das patentes mundiais situa-se exclusivamente em \textbf{G16H} (controle corporativo de prontuários reais) ou exclusivamente em \textbf{G09B} (plataformas genéricas de videoaulas). 

O diferencial tecnológico do SIGEA-GTT reside precisamente na \textbf{convergência sinérgica entre G16H e G09B}: a metodologia avançada de auditoria retrospectiva IHI-GTT é transposta e reengenheirada como uma ferramenta de simulação computacional ativa, provendo treinamento realista de alta fidelidade sem colocar em risco a integridade de pacientes reais e sem submeter os alunos ao estresse punitivo hospitalar.
